\subsection{Basic Encryption and Decryption}
To encrypt the file, I used the following command:

\begin{center}
    \begin{minipage}{0.8\linewidth}
        \begin{lstlisting}[language=bash, caption={Encrypting a file using AES-256-CBC}]
            openssl enc          \
                -aes-256-cbc     \
                -salt            \
                -pbkdf2          \
                -iter 10000      \
                -in message.txt  \
                -out message.enc
        \end{lstlisting}
    \end{minipage}
\end{center}

The following key was used when prompted:
\code{NcbSXXuJqXB6Tz1Xg3uAbUqFOwcfqG31yLEEs88S}.

To decrypt the file, the following command was used:

\begin{center}
    \begin{minipage}{0.8\linewidth}
        \begin{lstlisting}[language=bash, caption={Decrypting the file}]
            openssl enc          \
                -d               \
                -aes-256-cbc     \
                -pbkdf2          \
                -in message.enc  \
                -out message.dec
        \end{lstlisting}
    \end{minipage}
\end{center}

The same key was entered when prompted.

Below is a screenshot showing the execution of these commands:

\begin{minipage}{\linewidth}
    \centering
    \includegraphics[width=0.8\linewidth]{problem-8-step-1-1.png}
    \captionof{figure}{Task 1 -- Basic Encryption and Decryption -- Result}
\end{minipage}

The following image shows the comparison between the decrypted file
(\code{message.dec}) and the original (\code{message.txt}):

\begin{minipage}{\linewidth}
    \centering
    \includegraphics[width=0.8\linewidth]{problem-8-step-1-2.png}
    \captionof{figure}{Verification of file integrity using the \code{diff} command}
\end{minipage}
